\chapter{Agent Skills: the Operating Layer}\label{ch:skills}

The preceding chapters describe what the engine computes. This chapter describes the layer
through which the engine is most often \emph{operated}: a set of \emph{agent skills}, plain
Markdown instruction files that an AI coding agent loads when a request matches them, and that
turn a described task --- ``release particles near a point, backward, in three dimensions'',
``show where the decades differ'', ``why is this run slow'' --- into the concrete artefacts the
engine consumes: seeds, a configuration, job scripts, an analysis script, a figure script. The
skills contain no numerics of their own. They encode the order of operations and the checks that
past campaigns showed to be necessary, and they point to the documentation of
\cref{ch:arch,ch:input,ch:output,ch:performance} for everything else.

\section{What a skill is}\label{sec:skills-what}

A skill is a directory under \srcfile{.claude/skills/} whose entry point is a file
\srcfile{SKILL.md}. The file opens with a two-field header: a \emph{name} and a
\emph{description}. The description states only \emph{when} the skill applies --- the requests,
symptoms and boundary cases that should trigger it --- and deliberately does not summarise the
procedure, because an agent that finds the procedure in the description tends to act on the
summary and skip the body. The body is the procedure: an ordered workflow, the decisions that need
judgement with the condition that selects each option, and the checks that must hold. Longer
material --- templates, complete code recipes, measured tables, the history behind a rule --- lives
in files under \srcfile{references/} next to the skill, one level deep, and is read only when the
step that needs it names it. Two checks are shipped as runnable programs under \srcfile{scripts/}
(\cref{sec:skills-scripts}). An agent therefore reads a short entry file for every task and the
detail only for the task at hand.

\codeanchor{The seven skills are \srcfile{.claude/skills/kelp-run/}, \srcfile{kelp-config/},
\srcfile{kelp-experiment/}, \srcfile{kelp-analysis/}, \srcfile{kelp-visualize/},
\srcfile{kelp-diagnose/} and \srcfile{kelp-bench/}, each with a \srcfile{SKILL.md}; reference files in
their \srcfile{references/} subdirectories; the two check scripts in
\srcfile{kelp-run/scripts/}. The figure helpers the visualisation skill relies on are in
\srcfile{gallery/figure\_helpers.py}.}

\section{The seven skills}\label{sec:skills-seven}

\Cref{tab:skills} lists the skills with the request that triggers each and what it hands back.
They divide the work along the life of a study: setting up one run, writing its configuration,
organising a campaign of many jobs, reading the output, drawing it, diagnosing performance, and
benchmarking a machine.

\begin{table}[t]
\centering\footnotesize\setlength{\tabcolsep}{4pt}
\caption{The KELP agent skills. ``Hands back'' names the artefacts a skill produces; every skill
that prepares a computation stops before launching it (\cref{sec:skills-gate}).}
\label{tab:skills}
\begin{tabularx}{\linewidth}{@{}>{\raggedright\arraybackslash}p{29mm}>{\raggedright\arraybackslash}X>{\raggedright\arraybackslash}X@{}}
\toprule
Skill & Use it when & Hands back \\
\midrule
\code{kelp-run} & a tracking experiment is described (release location, direction, 2-D or 3-D,
diffusion, beaching, sampling, time window), or a dataset is used for the first time &
seed generator, configuration, verification commands and the run command; for a new dataset the
vertical-velocity sign and record-0 checks \\
\code{kelp-config} & a YAML configuration is written, edited, reviewed or debugged, or the engine
opened the wrong file, ignored a key or failed at startup & the configuration, a validation script
for generated or campaign configurations, a decision summary \\
\code{kelp-experiment} & an experiment is big (millions of particles, several processes or nodes),
long (beyond one job's wall clock), on a large mesh that may need cropping, or a speed and memory
showcase & a campaign folder with its rules file, report skeleton, ledger, config generator and
config check, job scripts, the resumable chain driver \\
\code{kelp-analysis} & anything is read or computed from trajectory output (maps, cohort fractions,
section crossings, run comparisons, reversibility, time-step ladders, costs from logs) & an
analysis script with the reading pattern, masks and denominators, and a plan whose numbers are
placeholders until the script prints them \\
\code{kelp-visualize} & a figure, animation or video is made from trajectory output or its caches,
or colours, palettes, labels or render speed need changing & a render script (extract, cache,
render), its commands, a caption file for each figure \\
\code{kelp-diagnose} & a run is slower than documented, silent, scales oddly with mesh size, costs
milliseconds per particle, or a proposed speed fix needs vetting & a diagnosis along a decision tree
and the measurement legs that decide it \\
\code{kelp-bench} & a machine is benchmarked or compared with the reference rows & one table row
from the self-contained single-node benchmark \\
\bottomrule
\end{tabularx}
\end{table}

The skills cross-reference one another rather than repeating each other
(\cref{fig:skills-routing}). \code{kelp-run} is the entry point for a single run and defers the
YAML mechanics to \code{kelp-config}; \code{kelp-experiment} sits above both and owns the
campaign order of operations; \code{kelp-analysis} owns the reading of trajectory output, which
\code{kelp-visualize} uses for its caches; \code{kelp-diagnose} is entered from the phase timings
that a verbose run prints and compares them with the measured rows of \srcfile{docs/performance.md}, and
\code{kelp-bench} produces one machine's row of the portable benchmark (\srcfile{benchmarks/single\_node/}).

\begin{figure}[t]
\centering
\begin{tikzpicture}[node distance=5mm and 7mm,
  sk/.style={rectangle,rounded corners=2pt,draw=lagblue,fill=lagblue!8,align=center,
             inner sep=3pt,minimum height=8mm,text width=29mm,font=\small},
  doc/.style={rectangle,rounded corners=2pt,draw=lagelem,fill=lagelem!8,align=center,
             inner sep=3pt,minimum height=7mm,text width=41mm,font=\footnotesize},
  req/.style={rectangle,rounded corners=2pt,draw=gray!60,fill=gray!6,align=center,
             inner sep=3pt,text width=112mm,font=\footnotesize},
  lk/.style={-{Stealth[length=4pt]},line width=0.7pt,draw=lagblue!70}]
  \node[req] (req) {a request: ``set up this run'' \;/\; ``plan this campaign'' \;/\; ``read this output''
    \;/\; ``draw this'' \;/\; ``why is it slow''\\[-1pt]
    \textit{matched against each skill's description (when to use), not its content}};
  \node[sk, below left=9mm and 26mm of req.south] (exp) {\code{kelp-experiment}\\\footnotesize campaign layer};
  \node[sk, below=of exp] (run) {\code{kelp-run}\\\footnotesize one run};
  \node[sk, below=of run] (cfg) {\code{kelp-config}\\\footnotesize the YAML};
  \node[sk, below=9mm of req.south] (ana) {\code{kelp-analysis}\\\footnotesize read output};
  \node[sk, below=of ana] (viz) {\code{kelp-visualize}\\\footnotesize figures, videos};
  \node[sk, below right=9mm and 26mm of req.south] (dia) {\code{kelp-diagnose}\\\footnotesize performance};
  \node[sk, below=of dia] (ben) {\code{kelp-bench}\\\footnotesize one machine};
  \node[doc, below=8mm of cfg] (d1) {\srcfile{docs/configuration.md}};
  \node[doc, below=8mm of viz] (d2) {\srcfile{gallery/} helpers};
  \node[doc, below=8mm of ben] (d4) {\srcfile{benchmarks/single\_node}};
  \node[doc, below=4mm of d4] (d3) {\srcfile{docs/performance.md}};
  \draw[lk] (req.south) -- ++(0,-3mm) -| (exp.north);
  \draw[lk] (req.south) -- (ana.north);
  \draw[lk] (req.south) -- ++(0,-3mm) -| (dia.north);
  \draw[lk] (exp) -- (run);  \draw[lk] (run) -- (cfg);
  \draw[lk] (ana) -- (viz);
  \draw[lk,dashed] (cfg) -- (d1); \draw[lk,dashed] (viz) -- (d2); \draw[lk,dashed] (ben) -- (d4);
  \draw[lk,dashed] (dia.east) -- ++(5mm,0) |- (d3.east);
  \draw[lk,dashed] (exp.east) to[out=0,in=180] (d3.west);
\end{tikzpicture}
\caption{How requests reach the skills. A request is matched against the skills' descriptions; the
selected \srcfile{SKILL.md} then points to the neighbouring skill it builds on (solid) and to the
documentation it defers to (dashed). Reference files under each skill's \srcfile{references/} are
opened only from the step that names them.}
\label{fig:skills-routing}
\end{figure}

\section{How the skill text is written}\label{sec:skills-design}

The skills are read by agents of different strength, and the way a rule is phrased decides whether
it is followed. The text therefore follows a small set of conventions.

\paragraph{Defaults and checks.} Each skill that produces a product opens with a short paragraph
that separates two kinds of guidance. \emph{Defaults} are the choices that worked in past work ---
a palette, a background, a layout, which tables to write, how an event is counted, which
uncertainty to quote. They are written as ``by default \dots; if the owner asks for something else,
do that'', because a default stated as a mandate makes a careful agent refuse what the owner
explicitly asked for. \emph{Checks} protect correctness whatever the style: the colour map's ends
agree with the label (an assertion in the script), dead particles are masked by status, lead times
are measured from each particle's own release, every number carries its denominator, cohorts the run
does not cover are dropped or censored, and a claim about a whole run is guarded by an assertion that
its last record is the run's end. Each check states its scope in the same sentence, so that
relaxing a default cannot silently relax the check. A handback ends with a section that lists every
default not used (with the owner's choice) and every check skipped (with the condition that did not
hold).

\paragraph{The launch gate.}\label{sec:skills-gate} The skills that prepare a run or a batch job (\code{kelp-run}, \code{kelp-config},
\code{kelp-experiment}) stop before starting it; the benchmark skill asks before running on a cluster
node. The prepared command is shown, and the agent asks for an explicit
confirmation of that command; a request that merely describes the run (``run it on 128 cores'') is
not a confirmation. The last line of their handback states that nothing was run.

\paragraph{Checks keyed to their condition.} A pre-launch check applies where its stated condition
holds: the vertical-velocity sign check for a three-dimensional run on a velocity not used before,
the record-0 check for a dataset not used before, the grid-hash cell check for a release in a
refined region of a non-uniform mesh. A check whose condition does not hold is skipped, with one
line saying why; the launch gate holds for every run.

\paragraph{Runnable code, no example numbers as answers.} Code blocks in the skills are complete
enough to run and were run before they were written down. Measured constants --- memory per particle and
throughput --- live in \srcfile{docs/performance.md} (a few per-event costs and the output sizing rule in
the skills' reference files, labelled as single measurements); a skill states the recipe that
derives the number for the task at hand and labels any quoted value as the measurement it came from.
A number from a worked example is never a substitute for the reader's own measurement.

\paragraph{Short entry files, reference files on demand.} Each \srcfile{SKILL.md} holds the routing,
the ordered workflow, the decision predicates and the checks; long templates and recipes live in
\srcfile{references/}. Two rules keep the split safe: the call line of every helper stays in the entry
file (a reader that never opens the reference still calls the helper correctly), and a reference that
depends on another links to it.

\section{The check scripts}\label{sec:skills-scripts}

Two pre-launch checks recur for every new dataset and are shipped as scripts, run from the clone
with the absolute path of the configuration.

\paragraph{Vertical-velocity sign.} \srcfile{w\_sign.py} loads the mesh and the dataset through the
configured reader with \conf{mode.w\_positive} forced to \code{up}, so that the negated internal
velocity is the file's own value (\cref{sec:vert-wsign}); it averages the file's vertical velocity
from \ang{2}\,S to \ang{2}\,N and from \ang{150}\,W to \ang{100}\,W, at the level nearest \SI{50}{m}, over records 1--30, and prints the mean with the sign
it implies. Equatorial upwelling makes the sign unambiguous there: a positive mean means the file
stores upward velocity. A store that does not cover this box stops with a request for another region
of known up- or downwelling. The configuration keeps \conf{w\_positive: auto} until the printed value
is written in, and the run script refuses \code{auto} for such a run.

\paragraph{Record 0.} \srcfile{record0.py} prints the fraction of exactly zero surface elevation at
the first, middle and last record. A model archive whose first record is the initial state of a cold
start shows a fraction near one there; the run then starts after the spin-up, which a periodic
forcing signal needs hours of records to develop --- skipping one record is not a spin-up.

\codeanchor{\srcfile{.claude/skills/kelp-run/scripts/w\_sign.py} and \srcfile{record0.py}; both use
\fn{load\_config}, \fn{get\_reader} and the reader's \fn{load\_mesh}, \fn{open\_dataset} and
\fn{read\_velocities}; \code{w\_sign.py} also uses \fn{unrotate\_lonlat} for a rotated mesh.}

\section{Skills and the documentation}\label{sec:skills-docs}

The skills are procedure; the documentation is reference. Option semantics, defaults and validation
rules are authoritative in \srcfile{docs/configuration.md}; the pipeline in
\srcfile{docs/architecture.md}; costs and sizing in \srcfile{docs/performance.md}; the restart
mechanics in \srcfile{docs/io\_and\_time.md}. A skill cites these rather than restating them, so that
a change in the engine is documented once. Where a skill and the documentation disagree, the
documentation and --- behind it --- the source take precedence, as throughout this manual.

\section{Maintaining a skill}\label{sec:skills-maintain}

A skill changes when a campaign shows that a step was missing, a rule was wrong, or a phrase led a
reader astray. Four practices keep the text reliable. Every code block in an edited skill is run
before it is committed. A changed rule is exercised by having a separate agent carry out a
representative task with the edited skill, and the artefact it produces --- not its summary --- is
reviewed, including a task in which the owner departs from a default on purpose, so that a new rule
does not override the owner. The review of an edit lists the errors the edit itself introduced, which
is where most new defects come from. Finally, the skills ship with the public engine: they carry
engineering practice and generic rules, and no study results.
